% Options for packages loaded elsewhere
\PassOptionsToPackage{unicode}{hyperref}
\PassOptionsToPackage{hyphens}{url}
\PassOptionsToPackage{dvipsnames,svgnames,x11names}{xcolor}
%
\documentclass[
]{article}

\usepackage{amsmath,amssymb}
\usepackage{iftex}
\ifPDFTeX
  \usepackage[T1]{fontenc}
  \usepackage[utf8]{inputenc}
  \usepackage{textcomp} % provide euro and other symbols
\else % if luatex or xetex
  \usepackage{unicode-math}
  \defaultfontfeatures{Scale=MatchLowercase}
  \defaultfontfeatures[\rmfamily]{Ligatures=TeX,Scale=1}
\fi
\usepackage{lmodern}
\ifPDFTeX\else  
    % xetex/luatex font selection
\fi
% Use upquote if available, for straight quotes in verbatim environments
\IfFileExists{upquote.sty}{\usepackage{upquote}}{}
\IfFileExists{microtype.sty}{% use microtype if available
  \usepackage[]{microtype}
  \UseMicrotypeSet[protrusion]{basicmath} % disable protrusion for tt fonts
}{}
\makeatletter
\@ifundefined{KOMAClassName}{% if non-KOMA class
  \IfFileExists{parskip.sty}{%
    \usepackage{parskip}
  }{% else
    \setlength{\parindent}{0pt}
    \setlength{\parskip}{6pt plus 2pt minus 1pt}}
}{% if KOMA class
  \KOMAoptions{parskip=half}}
\makeatother
\usepackage{xcolor}
\usepackage[top = 0.75in,bottom = 1in,left = 1in,right = 1in]{geometry}
\setlength{\emergencystretch}{3em} % prevent overfull lines
\setcounter{secnumdepth}{-\maxdimen} % remove section numbering
% Make \paragraph and \subparagraph free-standing
\makeatletter
\ifx\paragraph\undefined\else
  \let\oldparagraph\paragraph
  \renewcommand{\paragraph}{
    \@ifstar
      \xxxParagraphStar
      \xxxParagraphNoStar
  }
  \newcommand{\xxxParagraphStar}[1]{\oldparagraph*{#1}\mbox{}}
  \newcommand{\xxxParagraphNoStar}[1]{\oldparagraph{#1}\mbox{}}
\fi
\ifx\subparagraph\undefined\else
  \let\oldsubparagraph\subparagraph
  \renewcommand{\subparagraph}{
    \@ifstar
      \xxxSubParagraphStar
      \xxxSubParagraphNoStar
  }
  \newcommand{\xxxSubParagraphStar}[1]{\oldsubparagraph*{#1}\mbox{}}
  \newcommand{\xxxSubParagraphNoStar}[1]{\oldsubparagraph{#1}\mbox{}}
\fi
\makeatother


\providecommand{\tightlist}{%
  \setlength{\itemsep}{0pt}\setlength{\parskip}{0pt}}\usepackage{longtable,booktabs,array}
\usepackage{calc} % for calculating minipage widths
% Correct order of tables after \paragraph or \subparagraph
\usepackage{etoolbox}
\makeatletter
\patchcmd\longtable{\par}{\if@noskipsec\mbox{}\fi\par}{}{}
\makeatother
% Allow footnotes in longtable head/foot
\IfFileExists{footnotehyper.sty}{\usepackage{footnotehyper}}{\usepackage{footnote}}
\makesavenoteenv{longtable}
\usepackage{graphicx}
\makeatletter
\newsavebox\pandoc@box
\newcommand*\pandocbounded[1]{% scales image to fit in text height/width
  \sbox\pandoc@box{#1}%
  \Gscale@div\@tempa{\textheight}{\dimexpr\ht\pandoc@box+\dp\pandoc@box\relax}%
  \Gscale@div\@tempb{\linewidth}{\wd\pandoc@box}%
  \ifdim\@tempb\p@<\@tempa\p@\let\@tempa\@tempb\fi% select the smaller of both
  \ifdim\@tempa\p@<\p@\scalebox{\@tempa}{\usebox\pandoc@box}%
  \else\usebox{\pandoc@box}%
  \fi%
}
% Set default figure placement to htbp
\def\fps@figure{htbp}
\makeatother
% definitions for citeproc citations
\NewDocumentCommand\citeproctext{}{}
\NewDocumentCommand\citeproc{mm}{%
  \begingroup\def\citeproctext{#2}\cite{#1}\endgroup}
\makeatletter
 % allow citations to break across lines
 \let\@cite@ofmt\@firstofone
 % avoid brackets around text for \cite:
 \def\@biblabel#1{}
 \def\@cite#1#2{{#1\if@tempswa , #2\fi}}
\makeatother
\newlength{\cslhangindent}
\setlength{\cslhangindent}{1.5em}
\newlength{\csllabelwidth}
\setlength{\csllabelwidth}{3em}
\newenvironment{CSLReferences}[2] % #1 hanging-indent, #2 entry-spacing
 {\begin{list}{}{%
  \setlength{\itemindent}{0pt}
  \setlength{\leftmargin}{0pt}
  \setlength{\parsep}{0pt}
  % turn on hanging indent if param 1 is 1
  \ifodd #1
   \setlength{\leftmargin}{\cslhangindent}
   \setlength{\itemindent}{-1\cslhangindent}
  \fi
  % set entry spacing
  \setlength{\itemsep}{#2\baselineskip}}}
 {\end{list}}
\usepackage{calc}
\newcommand{\CSLBlock}[1]{\hfill\break\parbox[t]{\linewidth}{\strut\ignorespaces#1\strut}}
\newcommand{\CSLLeftMargin}[1]{\parbox[t]{\csllabelwidth}{\strut#1\strut}}
\newcommand{\CSLRightInline}[1]{\parbox[t]{\linewidth - \csllabelwidth}{\strut#1\strut}}
\newcommand{\CSLIndent}[1]{\hspace{\cslhangindent}#1}

\usepackage{setspace}
\onehalfspacing
\setlength{\parskip}{20pt}
\usepackage{orcidlink}
\usepackage{marvosym}
\makeatletter
\@ifpackageloaded{float}{}{\usepackage{float}}
\floatstyle{plain}
\@ifundefined{c@chapter}{\newfloat{excerpt}{h}{loexc}}{\newfloat{excerpt}{h}{loexc}[chapter]}
\floatname{excerpt}{Excerpt}
\newcommand*\listofexcerpts{\listof{excerpt}{List of Excerpts}}
\makeatother
\usepackage{xcolor}
\usepackage{marginnote}
\reversemarginpar
\newcommand{\paragraphnumber}[1]{\marginnote{\color{lightgray}\tiny{#1}}[0pt]}
\makeatletter
\@ifpackageloaded{caption}{}{\usepackage{caption}}
\AtBeginDocument{%
\ifdefined\contentsname
  \renewcommand*\contentsname{Table of contents}
\else
  \newcommand\contentsname{Table of contents}
\fi
\ifdefined\listfigurename
  \renewcommand*\listfigurename{List of Figures}
\else
  \newcommand\listfigurename{List of Figures}
\fi
\ifdefined\listtablename
  \renewcommand*\listtablename{List of Tables}
\else
  \newcommand\listtablename{List of Tables}
\fi
\ifdefined\figurename
  \renewcommand*\figurename{Figure}
\else
  \newcommand\figurename{Figure}
\fi
\ifdefined\tablename
  \renewcommand*\tablename{Table}
\else
  \newcommand\tablename{Table}
\fi
}
\@ifpackageloaded{float}{}{\usepackage{float}}
\floatstyle{ruled}
\@ifundefined{c@chapter}{\newfloat{codelisting}{h}{lop}}{\newfloat{codelisting}{h}{lop}[chapter]}
\floatname{codelisting}{Listing}
\newcommand*\listoflistings{\listof{codelisting}{List of Listings}}
\makeatother
\makeatletter
\makeatother
\makeatletter
\@ifpackageloaded{caption}{}{\usepackage{caption}}
\@ifpackageloaded{subcaption}{}{\usepackage{subcaption}}
\makeatother

\ifLuaTeX
\usepackage[bidi=basic]{babel}
\else
\usepackage[bidi=default]{babel}
\fi
\babelprovide[main,import]{canadian}
% get rid of language-specific shorthands (see #6817):
\let\LanguageShortHands\languageshorthands
\def\languageshorthands#1{}
\ifLuaTeX
  \usepackage[english]{selnolig} % disable illegal ligatures
\fi
\usepackage{bookmark}

\IfFileExists{xurl.sty}{\usepackage{xurl}}{} % add URL line breaks if available
\urlstyle{same} % disable monospaced font for URLs
\hypersetup{
  pdftitle={Data management is people management: on abstraction of data and labour in archaeological projects},
  pdfauthor={Zachary Batist},
  pdflang={en-CA},
  pdfkeywords={Information commons, collaborative commitments, divisions
of labour},
  colorlinks=true,
  linkcolor={blue},
  filecolor={Maroon},
  citecolor={Blue},
  urlcolor={Blue},
  pdfcreator={LaTeX via pandoc}}


\title{Data management is people management: on abstraction of data and
labour in archaeological projects}

\author{
      {Zachary
Batist \orcidlink{0000-0003-0435-508X} \href{mailto:zachary.batist@mcgill.ca}{\Letter}} \\
          McGill University
     \\
  }

\date{October 16, 2025}
\begin{document}
\maketitle
\begin{abstract}
Archaeological research is inherently collaborative, in that it involves
many people coming together to examine a material assemblage of mutual
interest by implementing a variety of tools and methods in sequence and
in tandem. Independent projects establish organizational structures and
information systems to help coordinate labour and pool information
derived thereof into communal data streams, which can then be applied
toward the production and publication of analytical findings. However,
these information commons are not egalitarian; project communities
establish practical expectations for how people in different positions
should engage with archaeological information based on their roles.
Contributing to and extracting from the commons is therefore scaffolded
by diverse yet converging commitments to a collective enterprise, which
are instilled through participation in a community of practice.

Through an abductive qualitative data analysis based on recorded
observations, interviews, and documents collected from three cases, this
paper highlights how digital systems designed to direct the flow of
information do so via the coordination of labour and the strategic
arrangement of human and object agency. Specifically, I highlight how
the information systems that archaeologists rely on to generate records,
internal reports, published papers and integrated datasets reify,
reinforce, and sometimes challenge entrenched power structures and
divisions of labour. For instance, I observe that each level of
documentation effectively communicates local understanding to a broader
collective by rendering prior work as series of formal processes
performed by interchangable agents; recognition of prior creative agency
therefore diminishes as narratives about the objects of interest are
further developed. At the same time, latent awareness of the practicial
circumstances, decisions and actions that contributed to records'
creation enable them to be recontextualized, as needed. However, this
ability to recontextualize data erodes at projects' boundaries, where
familiarity with creative circumstances fades.

By characterizing the documentary media that archaeologists use to
capture, organize and share knowledge as tools that facilitate
controlled communication, and which govern how certain people may
contribute to and access collectively-maintained knowledge, this paper
re-casts the formation of data commons as a social rather than a
technical enterprise. In highlighting some social and relational aspects
of data work that are often overlooked when developing local and global
research infrastructures, the paper puts forward a theoretical framework
that explains the awkwardness of using other people's data, and the
underwhelming results of initiatives to integrate data in contexts of
reuse.
\end{abstract}


\section{Introduction}\label{introduction}

Archaeological\paragraphnumber{[1]} projects are inherently
collaborative, in that they bring teams of people with a variety of
unique outlooks and experiences to examine material assemblages of
common interest. Archaeologists apply a multitude of tools and methods,
in sequence and in tandem and across different kinds of work settings,
to produce rich and heterogeneous data about their engagements with the
archaeological record. They adhere to professional norms and
expectations to organize the products of their collective labour, and
rely on a combination of digital and analog devices and protocols to
pool the outcomes of their respective tasks.

Advanced\paragraphnumber{[2]} computational systems have addressed many
technical barriers, but data integration is much more than just a
technical process; it also involves considerable collaborative labour to
facilitate commensurability of research outcomes and to ensure that all
parties effectively contribute to common goals. In this paper I draw
clear associations between project and data management, and their mutual
reliance on technological systems as means of maintaining control. As
such, it draws attention to the series of collaborative commitments and
broader social contexts from which we develop data management systems
and infrastructures.

Specifically,\paragraphnumber{[3]} the paper examines how
locally-circumscribed practices are re-framed as generic processes as
part of attempts to streamline the dissemination of data, and how this
permeates across project systems at various scales --- albeit unevenly
across various agents within project collectives.

\section{Methods and Data}\label{methods-and-data}

This\paragraphnumber{[4]} paper articulates behaviours and attitudes of
archaeologists who worked at two independent archaeological projects
which serve as cases. {[}about the cases{]}

It\paragraphnumber{[5]} is important to note that cases are not the
subjects of inquiry, and instead represent discrete instances that share
common reference to the overall research themes (Ragin 1992; Stake
2006). So while the cases are certainly not representative of the whole
archaeological discipline, they do enable me to reveal some
underappreciated systemic concerns that underlie data management. As
such, the findings are informed by the behaviours and attitudes
expressed to me by those who agreed to participate, and by my own
standpoint as a scholar of the culture and practice of science and of
the media and infrastructures that support data sharing, integration and
reuse.

Through\paragraphnumber{[6]} abductive qualitative analysis of
observations of archaeological practice, interviews with archaeologists
as they worked, and of documents they created, I captured diverse
experiences and outlooks that provide insight regarding how data
management infrastructures are valued and used in a variety of contexts.

The\paragraphnumber{[7]} findings presented here are derived from my
dissertation (Batist 2023), which produced three other papers deriving
from the same methods and data. See Batist and Roe (2024),
(\textbf{batist2025a?}), (\textbf{batist2025b?}) for further
methodological details.

\section{Findings}\label{findings}

Analysis\paragraphnumber{[8]} of archaeological materials and data is
dependent on systems of control. To control the flow of materials and
the structure of data, it is necessary to control people's behaviours.
In other words, the epistemic value that a project produces is
inherently dependent on the control that may be imposed on
archaeologists behaviours. Here I document some ways in which projects
scaffold archaeologists' behaviours through technical and administrative
means.

\subsection{Sampling}\label{sampling}

In\paragraphnumber{[9]} the cases I observed, specialist work often
involved specific sampling procedures that differed from, but also
occurred alongside, standard fieldwork practices. Specialists either
collected their own samples or instructed fieldworkers to collect
samples on their behalf and in ways that corresponded to their needs.
Trench supervisors were usually present to guide the specialist in their
trench, to help identify the locations from which specimens would likely
yield more useful results, and to be aware of changes to their work
environments that the sampling procedure may have caused. For example,
Marie, who was Case A's OSL specialist, was responsible for establishing
absolute dates for various stratigraphic units at the site, and was very
involved in the careful work of collecting samples from the
trenches.\textsuperscript{\hyperref[A73]{{[}A73{]}}} OSL dating
determines how long ago mineral grains were exposed to sunlight, and it
is critical to avoid exposing specimens to light. Sampling therefore had
to occur at night, and the involvement of trench supervisors was crucial
for determining where to draw the sample under such cumbersome
conditions.\textsuperscript{\hyperref[A96]{{[}A96{]}}} Moreover, it was
common practice to place a dosimeter where the sample was taken in order
to record background conditions that would have affected that specific
location, thus allowing calibration of the dose. Trench supervisors were
instructed to take care not to disturb these dosimeters as they
continued to excavate, until they would be collected at a later
time.\textsuperscript{\hyperref[A97]{{[}A97{]}}}

In\paragraphnumber{[10]} some cases, as when a specialist could not be
present during sample removal, it was deemed necessary to delegate
sampling work to a fieldworker who received training in basic sampling
procedures. This happened when a specialist had to leave the field
before the end of an excavation season, and but had enough foresight to
instruct a fieldworker to carry out the work on their
behalf.\textsuperscript{\hyperref[A73]{{[}A73{]}}} In this situation,
the fieldworker effectively worked in service of a specialist activity
and had to consider the criteria that would be most conducive to
effective specialist analysis.

A\paragraphnumber{[11]} similar phenomenon was observed at Case B, where
Chris was considered an expert in collecting spatial coordinates using
the differential global positioning system (DGPS) that forms the basis
of their geospatial analyses. Trench supervisors called upon Chris to
collect geospatial information from various targeted locations, and he
depended on their understanding of the trench and of the value of
specific spots that warranted data
collection.\textsuperscript{\hyperref[B17]{{[}B17{]}}} He also relied on
trench assistants to help him by holding a portable DGPS ``rover''
(which communicates with a spatially-calibrated base station located in
a fixed position), and talk among them tended to be about ensuring the
proper collection of a reliable sample of points (see
\textbf{?@fig-DGPS}).\textsuperscript{\hyperref[B18]{{[}B18{]}}} As
illustrated in Excerpt~\ref{exc-collecting-spatial-data}, Chris recorded
the geospatial point at the position where Oliver, the trench assistant,
placed the rover, while Liz, the trench supervisor, oversaw the process
and identified what needed to be
recorded.\textsuperscript{\hyperref[B18]{{[}B18{]}}}

\begin{excerpt}

\centering{

\begin{quote}
\textbf{Liz:} And then if you can just take one point there, and one
point right in the middle, umm that's all we need.\\
\textbf{Oliver:} Ok.\\
\textbf{Liz:} And then you're free to start digging and I'll help you
out in a minute.\\
\textbf{. . .}\\
\textbf{Chris:} So what are we doing?\\
\textbf{Oliver:} We're taking points.\\
\textbf{Liz:} Just one where you dug.\\
\textbf{Oliver:} Here?\\
\textbf{Liz:} Yep, yep. One in there, and then one over there is fine,
somewhere in the middle.\\
\textbf{. . .}\\
\textbf{Chris:} It's excavation unit {[}redacted identifier{]}?\\
\textbf{Liz:} Yeah. Closing.\\
\textbf{Chris:} Alright.??
\end{quote}

}

\caption{\label{exc-collecting-spatial-data}Communicating while
collecting spatial data.}

\end{excerpt}%

\subsection{Finds processing}\label{finds-processing}

Archaeological\paragraphnumber{[12]} projects uncover a lot of material
throughout the course of an excavation or survey, which must be sorted
and processed to facilitate subsequent analysis. Finds processing begins
in the field, where objects are first recovered and identified as finds
that will yield valuable insight into the lives of past people. At both
Cases A and B, fieldworkers divided these finds into broad-level
categories corresponding to materials, such as ceramic, bone and shell,
separating them into bags and buckets labelled with the excavation unit
in which they were found, as they emerged from the trenches (see
\textbf{?@fig-sieve}: E and F) and \textbf{?@fig-sorting}. At the same
time, they tossed stones that exhibit no anthropogenic qualities (see
\textbf{?@fig-sieve}: D). They also used sieves to separate loose
sediment from small objects, collecting those that they deem culturally
meaningful, and discarded the rest (see \textbf{?@fig-sieve}: A, B, C
and D).

Sometimes,\paragraphnumber{[13]} fieldworkers set aside materials that
they were unsure about and documented the findspot in more detail ---
``just in case'' --- before proceeding with their
work.\textsuperscript{\ref{A157}} This documentation was tentative, and
depending on whether a specialist would later determine that the find
was worth saving for analysis, may not have become part of the official
archaeological record.\textsuperscript{\ref{A156}, \ref{A157}}

It\paragraphnumber{[14]} is noteworthy that fieldworkers were expected
to classify finds on the basis of characteristics commonly recognized as
more natural, while study of their anthropogenic qualities was typically
reserved for finds analysis.\textsuperscript{\ref{A68}} In effect, a
distinction was made between fieldwork, which deals with
\emph{materials}, and finds analysis, which deals with \emph{artefacts}
and \emph{ecofacts}. Moreover, even when specialists worked in the
field, they operated as fieldworkers; their actions were constrained by
the practical circumstances of fieldwork and the position of fieldwork
within a broader apparatus of archaeological knowledge production. In
this sense, the division of labour was partly due to pragmatic concerns,
namely the fact that fieldworkers were usually occupied with their own
specific tasks that required mental concentration, especially when
instructed to do this work quickly.

After\paragraphnumber{[15]} hauling all the finds back to the dig house,
they needed to be cleaned by brushing the finds with
water.\textsuperscript{\ref{A69}, \ref{B9}} Trench assistants performed
this task after they have returned from the field, remaining productive
even during the late afternoons when it was too hot to perform strenuous
outdoor labour.\textsuperscript{\ref{B45}} In my unrecorded observations
I noticed that people talked about the day's work and their personal and
collective experiences, and speculated about the things they recovered
and which they were then revisiting in a secondary context. As
non-experts, their amateur analyses remain unexplored, but they still
felt an impetus to examine these objects and try to develop insights
from these experiences.

Cleaned\paragraphnumber{[16]} materials were then left out in the sun to
dry before being brought to a space dedicated for finds analysis and
storage.\textsuperscript{\ref{B46}} The materials were dumped on a broad
table, and specialist's assistants quickly sorted through them. They
discarded any non-anthropogenic materials collected in error, then
counted and weighed all remaining artefacts, before performing further
analysis based on their typological and technological
characteristics.\textsuperscript{\ref{A158}}

At\paragraphnumber{[17]} Case A, I observed how Jolene, a lithics
specialist, went through and sorted the material on the table, sometimes
mumbling her reasoning aloud.\textsuperscript{\ref{A161}} After she
settled on her arrangement, she dictated her description of the
assemblage as a whole, in a fragmented tone that insinuates strong
intention behind each word, and Maude, her assistant, wrote down this
speech.\textsuperscript{\ref{A161}} Maude then recorded the quantities
of artefacts in each category, as laid out by Jolene, and entered this
information into a spreadsheet.\textsuperscript{\ref{A161}} Whereas
Jolene was the source of creative knowledge production, Maude was
responsible for embedding these insights within the project's
organizational apparatus.\textsuperscript{\ref{A162}}

Maude\paragraphnumber{[18]} also maintained a checklist to track all
aspects of work performed on the materials, including information about
the materials' provenance, who processed the materials, when the work
occurred, and detailed parameters that structured processing
tasks.\textsuperscript{\ref{A160}} This checklist resembled the records
maintained by the project's database manager to trace what records await
integration into the database and to ensure that all information was
properly recorded.\textsuperscript{\ref{A87}, \ref{A163}}. These records
existed to track progress and verify the integrity of work. No aspect of
finds processing involved maintaining a journal or recording substantial
interpretations of the materials, which revealed a failure to recognize
that any meaningful insights could arise during these stages of work.
Finds processing was therefore positioned as a means to an end, that
being the arrival of cleaned and sorted finds upon the specialists'
desks, where meaningful archaeological knowledge would be
produced.\textsuperscript{\ref{A129}, \ref{A69}, \ref{A41}, \ref{B9}}

\subsection{Integrating and aligning
perspectives}\label{integrating-and-aligning-perspectives}

After\paragraphnumber{[19]} performing their work, specialists reported
back to the project with their findings. This involved providing the
project with documents that contained and explained the analytical
results. These documents typically comprised electronic spreadsheets and
typeset documents, but sometimes also included lab journals, paper
recording sheets, source code, and digital images. These media varied in
their functional roles and value.

For\paragraphnumber{[20]} instance, the spreadsheets that archaeological
specialists produce in the cases I examined were the formally organized
results of systematic analysis. They were organized on a
sample-by-sample basis, associated with discrete units from which each
sample element was taken (see
\textbf{?@fig-archaeometry-spreadsheets}).\textsuperscript{\hyperref[A6]{{[}A6{]}},~\hyperref[A98]{{[}A98{]}}}
These records were analytical results in that they were the products of
analytical processes, but they were also data in the sense that they
were considered to be grounded and stable records of systematic
observation that may serve as a sound basis for further investigation.

Specialist\paragraphnumber{[21]} reports, on the other hand, were
summaries of a specialist's findings. They described overall trends that
specialists observed in their preliminary analysis of the findings,
highlighted their potential significance, and proposed ways to
ameliorate the results or relate them to other aspects of the
project.\textsuperscript{\hyperref[B11]{{[}B11{]}}} In other words,
specialist reports contextualized the analysis and related the findings
to the project's general research objectives. They presented conclusions
that aligned or conflicted with alternative perspectives, settled or
provided greater confidence to particular claims, and suggested aspects
of work that warranted further attention. They did these things in an
embodied manner, and in ways that reflected the pragmatic situation of
their work; they served to contextualize the data presented in
spreadsheets, which lacked any such capability on their own.

Specialist\paragraphnumber{[22]} datasets were integrated into projects'
relational databases by importing the relevant records as their own sets
of tables, with explicit links drawn between them and the database
backbone, which reflected the organizational structure of archaeological
entities encountered in the field (see
\textbf{?@fig-archaeometry-spreadsheets}).\textsuperscript{\hyperref[A98]{{[}A98{]}}}
For example, samples that were drawn from a particular excavation unit
were typically listed in a specialist table, but were linked to an
overarching index of all excavation
units.\textsuperscript{\hyperref[A6]{{[}A6{]}}} Analytical results could
therefore be retrieved for particular parts of a site, alongside
findings derived from different kinds of analysis performed on material
deriving from the same contexts.

The\paragraphnumber{[23]} material that was physically distributed and
sent out for specialized analysis was thus conceptually re-integrated in
informational terms. Diverse kinds of material were rarely examined in
close physical proximity outside of fieldwork, though the potential to
integrate abstract measurements about them was highly
valued.\textsuperscript{\hyperref[B19]{{[}B19{]}},~\hyperref[C17]{{[}C17{]}},~\hyperref[C18]{{[}C18{]}},~\hyperref[C19]{{[}C19{]}}}
However, this actually rarely happened in ways that leveraged the
computational potential of these efforts. Specialists provided their own
preliminary analyses in written and oral reports, in which they examined
the materials under their own
purview.\textsuperscript{\hyperref[B1]{{[}B1{]}},~\hyperref[B11]{{[}B11{]}}}
Analysis was therefore siloed, and findings deriving from these analyses
were subsequently
integrated.\textsuperscript{\hyperref[A99]{{[}A99{]}},~\hyperref[A89]{{[}A89{]}},~\hyperref[B12]{{[}B12{]}}}

While\paragraphnumber{[24]} specialists' interests may have overlapped
with those of other members of the project, directors preferred to have
one specialist for each means of investigation. At the same time, some
degree of overlap among specialists who examine the same kinds of
materials from different methodological perspectives was expected and
encouraged. For example, the geoarchaeologists at Case A shared a strong
affinity based on their mutual work with geological materials, despite
the fact that they specialized in different geological
processes.\textsuperscript{\hyperref[A85]{{[}A85{]}},~\hyperref[A95]{{[}A95{]}}}
This benefited the project immensely, since it allowed each specialist
to reflect upon the limitations of their approach and to harmonize that
approach with those of the other specialists.

At\paragraphnumber{[25]} the same time, however, the value of these
complementary perspectives has yet to be demonstrated when presenting
findings to the research community. Much of the discussion among
geoarchaeologists at Case A, for instance, concerned strategic planning,
including planning where to excavate, prioritizing where to take
samples, and considering the degree of confidence that should be placed
in certain findings. In other words, these were process-oriented
discussions that reflected on decision-making processes. This
information was not effectively represented in the database and in
reports, which were more meant to articulate research outputs in a clear
and confident manner. Inputting analytical findings into a database was
therefore a lossy process that failed to capture much of the pragmatic
knowledge and critical reflection provided by individuals and
collectives.

\subsection{Informal consultation}\label{informal-consultation}

Additional\paragraphnumber{[26]} interactions between archaeologists and
external members of their research communities also occurred through
consultation with key figures, whom projects called in to either provide
informed insights as a prelude to broader dissemination of a project's
findings or to provide expert opinions on how work should proceed.
Consultants were specialists in their own right, but their involvement
was usually limited to one-off interactions that precluded long-term
commitment to the project.

As\paragraphnumber{[27]} Case A, a consultant was called in to examine a
peculiar and anomalous material configuration that confounded project
participants. As illustrated in Excerpt~\ref{exc-consulting-dialogue},
Basil asked Andreas, a micromorphology specialist, to perform a
preliminary, \emph{ad-hoc} examination of the microstratigraphy in a
trench that exhibited evidence of hearth
activity:\textsuperscript{\hyperref[A118]{{[}A118{]}},~\hyperref[A119]{{[}A119{]}},~\hyperref[A120]{{[}A120{]}},~\hyperref[A121]{{[}A121{]}}}

\begin{excerpt}

\centering{

\begin{quote}
\textbf{Andreas:} And on this side, on this side, is the, the fire base,
sand, sand right below the, the hearth. \dots This sand is de-calcified,
which again implies kind of, {[}left on{]} the surface in order to
de-calcify. And on this surface, the guys came and put fire. Actually
the matrix of the black stuff is sand again. \dots So the whole thing is
kind of in play but not in sequence. \dots Of course we don't see any
more ashes because they have been dissolved away. Actually what you see
under the microscope is just sand grains surrounded by {[}unclear{]}
made of very fine charcoaled dust, dust charcoal. That's why you don't
find, I guess, charcoal inside, probably {[}a{]} few bits.\\
\textbf{Basil:} Yeah, it's nothing. Yeah it's a very very tiny amount.\\
\textbf{Andreas:} Yeah, I mean okay, it's like putting a fire down on
the sand. You know, the sand is something that moves, always, around.
Even people when we're trampling it, you know animals, you know \dots.
Sadly, nothing would survive, \dots so much not in a \dots but would end
up in dust, something like that.\\
\textbf{Alfred:} Oh well.\\
\textbf{Andreas:} But still, it's {[}held up{]} in place, and made on
the, on the material that is just below. So they are, let's say, kind of
the same, umm, story.\\
\textbf{Basil:} Excellent.\\
\textbf{Alfred:} That's good. Yeah.\\
\textbf{Andreas:} I mean you can date, \dots The accumulation of the
sand below and the fireplaces on top, they might be about the same age,
more or less. I mean, if you are lucky and you \dots you can separate
these two events, it would be nice but I think it's, it wouldn't be
very\ldots{}\\
\textbf{Alfred:} {[} unclear, but something about carbonized remains
{]}\\
\textbf{Basil:} We have a few carbonized remains. And we, we just took
another 167 litres of flotation so our, we have our archaeobotanist
coming through next week.\\
\textbf{Andreas:} Yeah, then of course. \dots It is, there's no doubt
about it.\\
\textbf{Alfred:} Yeah.\\
\textbf{Andreas:} Actually yeah, ok I forgot. There's quite a few burnt
bone inside but are reduced in the size of sand. Let's say less than a
half a millimetre, but they are burnt. Again, because you know, they are
so {[}can't hear because of the sieve nearby{]}, they've degraded in
that small sizes.\\
\textbf{Basil:} Well, that's the only bone we have from the site so
far.\\
\textbf{Andreas:} I can show you under the mic, I mean they are tiny
like \dots sand grains. And they are burned too.\\
\textbf{Basil:} I wonder if we could do anything further with that.\\
\textbf{Andreas:} Other than dating, I don't think, what else? Yeah.
Yeah, I think more or less, you can't do more.\\
\textbf{Basil:} So what---\\
\textbf{Andreas:} Ahh, I mean, ok, in theory, some more fancy
biochemical analysis? I don't know, I don't know if these things have
survived. But why not? There have been, lately, some umm biochemical
analyses on these kinds of stuff.\textsuperscript{\ref{A121}}
\end{quote}

}

\caption{\label{exc-consulting-dialogue}Andreas provides guidance on how
to collect meaningful information.}

\end{excerpt}%

In\paragraphnumber{[28]} this exchange, Andreas remarked that there was
enormous potential to apply more in-depth micromorphological and
instrumental analyses in this particular trench. He provided guidance on
how to expand excavation around this trench on the basis of his
understanding of the geological processes at work on the site and how
they would affect the archaeological traces of hearths, and he suggested
new forms of analysis that the director had not yet considered. Each
participant played a particular role here: Andreas provided his brief
analysis as a visiting and non-committed specialist; Alfred, who was the
project's dedicated micromorphology specialist, recalled specific
aspects that can be leveraged to support his own investigations; and
Basil used his authority as project director to decide that these ideas
were worth pursuing in greater
depth.\textsuperscript{\hyperref[A120]{{[}A120{]}},~\hyperref[A121]{{[}A121{]}}}

Similarly,\paragraphnumber{[29]} since Case B did not employ a full-time
zooarchaeologist, Chris (who was also one of the co-directors)
occasionally sent pictures of faunal remains they recovered to a
colleague who would quickly identify them. For example, I observed one
instance when Chris was fairly confident that a long bone they found was
a cow's femur but he wanted to verify that it was not a human bone,
which would have called for very different excavation procedures and
project management decisions.\textsuperscript{\hyperref[B31]{{[}B31{]}}}

Additionally,\paragraphnumber{[30]} I observed that projects sometimes
seek to leverage consultants' authority on a subject to bolster
confidence in cases when their findings rested on shaky
foundations.\textsuperscript{{[}A122{]},{[}A123{]}{[}A124{]}} This was
exemplified by the way that Case A handled its anxieties regarding the
extremely colluvial nature of the site. Colluviation complicates any
effort to date stratigraphic units and had resulted in many artefacts
having extremely weathered surfaces that make them more difficult to
characterize. Despite his confidence in their claims due to his daily
involvement examining these artefacts, and his acknowledgement that the
stratigraphy is imperfect, Basil knew that he would have a hard time
convincing others of the nature of the things that they found. He
therefore involved Denise, a major figure in lithics analysis who was
responsible for publishing the most definitive assemblages of the kind
of materials recovered at Case A, to examine their materials, as
articulated in
Excerpt~\ref{exc-warrant-for-consultantation}:\textsuperscript{{[}A124{]},{[}A125{]}}

\begin{excerpt}

\centering{

\begin{quote}
\textbf{Basil:} \ldots{} for decades, there have been accepted means of
scientifically representing the archaeological record, umm to act as
evidential bases for your claims. It's understood, that you know, it's
impractical for every single interested archaeologist to come to see
your site, to see your material. You could be in Australia and the site
could be in Peru, umm and so, without that tactile, immediate
relationship with the evidence basis that somebody's uhh excavated,
there are ways of representing it through photography and line
illustration. Umm, that always used to seem to be enough, but it seems
like again, it comes down to this contentious nature of our site. The
fact that we are, we allegedly have such an early site where such early
sites are not meant to be, I'm starting to discover that some of what
used to be the accepted means of, accepted ehh evidential umm bases,
don't seem to be good enough for everyone. And so it's been very
important to have people like Denise come and see this stuff in the
flesh. Yes, she's seen the drawing, yes, she's seen the photograph, but
for her to see it first hand, handle it, pick it up, query it, umm look
at, look at assemblages, as opposed to cherry picked illustrated pieces,
has been fundamental to convince her that we really have what we say we
have.{[}A124{]}
\end{quote}

}

\caption{\label{exc-warrant-for-consultantation}Basil explains his
rationale for inviting Denise to examine the archaeological material.}

\end{excerpt}%

At\paragraphnumber{[31]} the same time, Denise informed Basil about some
of her concerns with the material, and in
Excerpt~\ref{exc-result-of-consultation} Basil recalls how this
constructive criticism provided him with an understanding of how to
proceed to accommodate these
issues.\textsuperscript{{[}A125{]},{[}A126{]}}

\begin{excerpt}

\centering{

\begin{quote}
\textbf{Basil:} So the superstar material that Lauren was excavating,
that we thought was Levallois blade cores, she disagreed. She thinks the
material is fabulous, she's never seen this material before in Greece,
she thinks it's Aurignation and looks much more like what you would find
in France. She's blown away. She's blown away by---\\
\textbf{Zack:} But it's still Middle Pal?\\
\textbf{Basil:} Very early Middle Pal. But, we had shown her our
superstar pieces. This is where it---\\
\textbf{Zack:} I was in the trench that first day when we, and I was
like cores, wow, tons of them.\\
\textbf{Basil:} I mean that stuff is--- but we had shown her some of our
superstar pieces from the survey, like here's a Levallois point, classic
thing. But she's like hmm, it's just one piece. Haha you know, it's just
like, she's not exactly {[}unclear{]}, but she has very high
standards.\\
\textbf{Zack:} So is she sticking with, is she saying that Middle Pal
was a thing though?\\
\textbf{Basil:} She was like, so she was looking at these, these little
bits and pieces, and like, well, the stuff, you know, we showed her a
whole bunch of stuff, it's like, we think this is Levallois, she's like
no, it's this. And then we show her some stuff from the survey and it's
like, hmm it's one piece and it could be, but it's just one piece,
that's not really enough. And then finally we showed her the material,
Jolene showed her material from {[}the trench{]}, Alice's new trench,
and she's like yes, this is Middle
Palaeolithic.\textsuperscript{\ref{A125}}
\end{quote}

}

\caption{\label{exc-result-of-consultation}Basil explains the feedback
he received from Denise, and what this means for future work.}

\end{excerpt}%

From\paragraphnumber{[32]} these examples we see that members of the
archaeological community who were not active project participants
sometimes supplemented the input provided by specialist project
participants. As people who were external to the project, they were not
committed to following through on their advice; they instead relayed
their interpretations to specialists who are already members of the
project and who were capable of implementing their
guidance.\textsuperscript{{[}A127{]},{[}A128{]}} Moreover, as outsiders,
they were not forced to adhere to the systematic recording systems that
scaffolded internal procedures, and the informal nature of their
contributions was in fact a crucial aspect that imbued them with value.

\section{Discussion and Conclusion}\label{discussion-and-conclusion}

Projects\paragraphnumber{[33]} experience immense pressures to render
observations in ways that are amenable to formal and functional
analysis, which in turn mandate rigid and decisive data collection
processes to instill confidence in the stability and evidentiary nature
of archaeological data.\textsuperscript{\ref{A176}, \ref{A89}} As such,
particular kind of analytic process comes to dictate the conduct of
other aspects of work, and involves reducing the interpretive agency of
actors involved in all aspects of the data's constitution and
processing.

This\paragraphnumber{[34]} results in prioritization of archaeological
materials as formally structured and abstract entities, devoid of
character in their own right and always perceived as specific instances
of a more general class. The only person who can recognize a thing as
meaningful is someone who is deemed capable of granting it meaning. This
is exemplified by my observations of work performed by trench
assistants, who essentially serve as unskilled labourers who are
shuffled around a project to perform menial tasks that, almost
exclusively, concern materials that have not yet been assigned
archaeological character (\textbf{batist2025?}). Beyond operating as
sensory instruments, fieldworkers serve others' creative efforts by
cleaning artefacts and performing preliminary sorting tasks. Those
involved in this type of work prepare materials for analysis that others
will perform, namely specialists. Similarly, trench assistants move
sediment and rocks out of a trench, but it is up to the supervisor, who
is responsible for recording the significance of these materials, to
make meaning out of them. Trench assistants are cast as mere sensors,
whose identifications of a thing, if taken seriously at all, are based
on the character of shapes and textures they encounter, and not on any
underlying archaeological theory or model.

Their\paragraphnumber{[35]} interchangeability, owing to the limited
expectations made of them (simply responding to prompts), renders
fieldworkers invisible and removes any potential for imagining data
collection as discursive and situated. This supports the absorption of
fieldworkers into distributed data-recording system, thereby reducing
the impacts of their situated perspectives and presenting their work as
a form of objective sensory understanding.

At\paragraphnumber{[36]} the same time, interpretation is delegated to
others who are not as present at the site, and whose distance from the
archaeological encounter provides a kind of freedom to take the data at
face value, rather than to grapple with their discursive origins. This
contributes to a sense that so-called ``data-driven'' analysis generates
new knowledge synthesized from a series of objective facts derived from
generic, predictable, and protocol-driven sensing devices (Huggett
2022). Archaeological information infrastructures are thus designed to
enable the appropriation of direct observations by those with the means
of extracting more meaningful value from them.

Moreover,\paragraphnumber{[37]} the paper observed how there is greater
flexibility for those higher up in the chain to explain the context of
the work they perform. However, the advent of open data sharing,
integration and reuse adds another layer of abstraction, whereby entire
projects and the totality of the data they collect are rendered as
generic instances of some greater class. The discomfort that project
directors sense upon having their creative agency reduced is akin to
that which was experienced during the advent of tighter control over
fieldwork recording systems (as documented by Chadwick (1998) and Thorpe
(2012)), and which we now take for granted.

\begin{center}\rule{0.5\linewidth}{0.5pt}\end{center}

This\paragraphnumber{[38]} paper frames the technologies that
archaeologists use to re-situate individual archaeologists in relation
to project collectives as means through which managerial protocols are
carried out. Specifically, it documents a few underappreciated aspects
of archaeological projects' sociotechnical arrangements that should be
accounted for more thoroughly:

\begin{enumerate}
\def\labelenumi{\arabic{enumi}.}
\tightlist
\item
  Systems designed to direct the flow of information do so via the
  coordination of labour, and the strategic arrangement of human and
  object agency.
\item
  These systems are \emph{attempts} at control, which are actively
  resisted by those who work under them, and passively resisted by the
  wild nature of the things that that they attempt to tame.
\item
  Informal and socially-mediated forms of scholarly communication are
  de-prioritized and excluded from data streams ascribed with greater
  scholarly legitimacy, despite the fact that this information is
  crucial for developing a sense of mutual understanding that underpins
  comprehension of formal datasets.
\item
  Boundaries, whether they restrict membership within a collective,
  delimit a project's scope, or limit the time frame under which a
  project operates, have practical positive value, and are not just
  arbitrary impediments.
\item
  Information systems and the institutional structures that support them
  tend to reinforce and reify existing power structures and divisions of
  labour, including implicit rules that govern ownership and control
  over research materials and that designate who may benefit from their
  use.
\end{enumerate}

\section{References}\label{references}

\phantomsection\label{refs}
\begin{CSLReferences}{1}{0}
\bibitem[\citeproctext]{ref-batist2023a}
Batist, Zachary. 2023. {``Archaeological Data Work as Continuous and
Collaborative Practice.''} PhD thesis, Toronto: University of Toronto.
\url{https://hdl.handle.net/1807/130306}.

\bibitem[\citeproctext]{ref-batist2024}
Batist, Zachary, and Joe Roe. 2024. {``Open {Archaeology}, {Open
Source}? {Collaborative} Practices in an Emerging Community of
Archaeological Software Engineers.''} \emph{Internet Archaeology}, no.
67 (July). \url{https://doi.org/10.11141/ia.67.13}.

\bibitem[\citeproctext]{ref-chadwick1998}
Chadwick, Adrian. 1998. {``Archaeology at the Edge of Chaos: Further
Towards Reflexive Excavation Methodologies.''} \emph{Assemblage}
3:97--117.

\bibitem[\citeproctext]{ref-huggett2022}
Huggett, Jeremy. 2022. {``Is {Less More}? {Slow Data} and {Datafication}
in {Archaeology}.''} In \emph{Critical {Archaeology} in the {Digital
Age}: {Proceedings} of the 12th {IEMA Visiting Scholar}'s {Conference}},
edited by Kevin Garstki, 97--110. Los Angeles, CA: Cotsen Institute of
Archaeology Press.
\url{https://escholarship.org/uc/item/0vh9t9jq\#page=112}.

\bibitem[\citeproctext]{ref-ragin1992}
Ragin, Charles C. 1992. {``Casing and the Process of Social Research.''}
In \emph{What Is a {Case}? {Exploring} the {Foundations} of {Social
Inquiry}}, edited by Charles C. Ragin and Howard S. Becker, 217--26.
Cambridge University Press New York.

\bibitem[\citeproctext]{ref-stake2006}
Stake, Robert E. 2006. \emph{Multiple Case Study Analysis}. New York:
Guilford Press.

\bibitem[\citeproctext]{ref-thorpe2012}
Thorpe, Reuben. 2012. {``Often {Fun}, {Usually Messy}: {Fieldwork},
{Recording} and {Higher Orders} of {Things}.''} In \emph{Reconsidering
{Archaeological Fieldwork}: {Exploring On-Site Relationships Between
Theory} and {Practice}}, edited by Hannah Cobb, Harris Harris Oliver J.
T., Cara Jones, and Philip Richardson, 31--52. Springer, Boston, MA.
\url{https://doi.org/10.1007/978-1-4614-2338-6_3}.

\end{CSLReferences}

\subsection{Acknowledgments}\label{acknowledgments}

I\paragraphnumber{[39]} extend warm thanks to Costis Dallas, Matt Ratto,
Ted Banning, Jeremy Huggett and Ed Swenson for supervising my work and
providing critical feedback as I conducted my doctoral research, which
this paper is based upon. I am also very grateful to the anonymous
reviewers for their constructive evaluation. Of course, this work would
not have been possible without the anonymous research participants who
allowed me to observe and interview them as they worked, and to
articulate their actions and outlooks.

\subsection{Funding}\label{funding}

This\paragraphnumber{[40]} work is derived from the author's doctoral
dissertation, \emph{Archaeological data work as continuous and
collaborative practice} (Batist 2023), which was supported by the
Canadian Social Sciences and Humanities Research Council Doctoral
Fellowship (Award ID: 752-2019-2233).

\subsection{Competing Interests}\label{competing-interests}

The\paragraphnumber{[41]} author states no conflicts of interest.

\subsection{Author Contributions}\label{author-contributions}

Zachary\paragraphnumber{[42]} Batist is the sole author of this work. He
defined the scope of the study and identified suitable cases for
inclusion, collected and processed all data, performed analysis,
interpreted the findings, created all the figures, and wrote the paper.

\subsection{Informed Consent}\label{informed-consent}

Informed\paragraphnumber{[43]} consent has been obtained from all
individuals included in this study, in compliance with the University of
Toronto's Social Sciences, Humanities, and Education Research Ethics
Board, Protocol 34526.

\subsection{Data Availability
Statement}\label{data-availability-statement}

The\paragraphnumber{[44]} data generated and analyzed during the current
study are included in this published article's supplementary files.

\subsection{Author Bios}\label{author-bios}

Zachary\paragraphnumber{[45]} Batist obtained his PhD from the
University of Toronto's Faculty of Information. His research explores
the collaborative commitments inherent throughout archaeological
practice, especially relating to data management and the constitution of
information commons. He currently works as a Postdoctoral Researcher at
the Department of Epidemiology, Biostatistics and Occupational Health in
the School of Public and Global Health at McGill University, where he
investigates the collaborative, technical and administrative structures
that scaffold data harmonization in epidemiological research.

\subsection{Figure Captions}\label{figure-captions}

\section{Supplementary Materials}\label{supplementary-materials}




\end{document}
